\chapter{Force Estimation Quality}
\label{ch:evaluation}

This chapter defines what we mean by \emph{good force estimation} in the
context of linear-mapping compliance, specifies a single benchmark protocol
that every trained checkpoint is evaluated on, and reads off the effect of
each architectural choice from the resulting table. Because the downstream
compliance map is linear ($\vec{v}_{\text{cmd}} = k\,\hat{\vec{f}}$), the
walker's compliant behaviour is a deterministic function of the estimator
output and a tuned scalar. Estimator quality is therefore the only
research object, and the chapter is organised around it.

%==============================================================================
\section{What is good force estimation?}
\label{sec:eval_definitions}

We report six measurements for every checkpoint. No measurement is given a
pass/fail threshold; they are descriptive and read together.

\begin{table}[h]
\centering
\caption{The six measurements of estimator quality used throughout the chapter.}
\label{tab:eval_metrics}
\begin{tabular}{ll}
\hline
Measurement & Definition \\ \hline \hline
Magnitude MAE & $\operatorname{mean}\left(\|\hat{\vec{f}}-\vec{f}^{*}\|\right)$ over samples with $\|\vec{f}^{*}\|\ge 1\,$N \\
Relative MAE & MAE normalised by $\operatorname{mean}(\|\vec{f}^{*}\|)$ \\
Angular error (xy) & $\operatorname{median}\left(\angle(\hat{\vec{f}}_{xy},\vec{f}^{*}_{xy})\right)$, in degrees \\
Magnitude-conditional MAE & MAE stratified by $\|\vec{f}^{*}\|$ into fixed bins \\
Per-direction MAE & MAE as a function of the applied force azimuth \\
Responsiveness & 90\%-rise time of $\hat{\vec{f}}$ in response to a force step \\
Zero-force behaviour & distribution of $\|\hat{\vec{f}}\|$ when $\|\vec{f}^{*}\|=0$ \\
\hline
\end{tabular}
\end{table}

The two views that carry most of the story are the \emph{magnitude-conditional
MAE} (\cref{sec:eval_training_distribution}) and the \emph{zero-force
behaviour} under joystick driving (\cref{sec:eval_self_motion_fps}). The
first exposes whether an estimator generalises across the force spectrum;
the second exposes whether it hallucinates forces from its own motion.

Because the linear mapping is deterministic, we report walker-side
quantities only as sanity checks: (i) the tracking-reward delta of the
compliant walker relative to the base walker under zero applied force, and
(ii) the consistency of $\vec{v}_{\text{cmd}}$ across applied magnitudes at
fixed direction. Both fall out of the static sweep protocol below without
any additional experiments.

%==============================================================================
\section{Benchmark protocol}
\label{sec:eval_protocol}

Every checkpoint is evaluated with the same four-test protocol. Each test
corresponds to a single evaluation script
(\cref{sec:eval_run_guide}) and produces a fixed set of output files that
populate the tables in this chapter.

\subsection{Static 360\textdegree{} force sweep, standing}
The robot stands on flat terrain. External forces are applied with
$\|\vec{F}\|\in\{5, 10, 20, 30, 40, 50\}\,$N at $12$ evenly-spaced
azimuths, each held for \unit[4]{s}, $N_{\text{trials}}=20$. Populates
the magnitude, directional, magnitude-conditional, and per-direction
metrics of \cref{tab:eval_metrics}.

\subsection{Step-response}
The robot stands on flat terrain. At $t=0$ a constant force is applied
instantaneously at one of four azimuths $\{0,90,180,270\}$\textdegree{} and
held for \unit[2]{s}. Populates the responsiveness metric.

\subsection{Zero-force standing (flat and slope)}
The robot stands on (a)~flat ground and (b)~a \unit[15]{deg} slope for
\unit[10]{s} each, with no external force. Populates the flat and sloped
cells of the zero-force row of the master table.

\subsection{Zero-force joystick trajectory}
\label{sec:eval_protocol_joystick}
The robot is driven by a scripted command trajectory consisting of
forward/backward walking, turning in place, lateral strafing, and
standing, \unit[30]{s} per segment. No external force is applied at any
point. This is the test introduced specifically for the self-motion
false-positive question of \cref{sec:eval_self_motion_fps}.

\medskip

The protocol is executed in simulation for every checkpoint in
\cref{sec:eval_runs}. Tests~1, 3 and 4 are additionally executed on the
real robot for the subset identified in \cref{sec:eval_real}.

%==============================================================================
\section{Runs included in the study}
\label{sec:eval_runs}

Three groups of checkpoints are considered. All three use the same
locomotion environment and differ only in the aspects listed.

\begin{table}[h]
\centering
\caption{Run roster. All runs use the same low-level locomotion environment.
Max-force is the upper bound of the per-axis force distribution used during
training; mean $\|\vec{F}\|$ is the empirical training-time magnitude.}
\label{tab:eval_runs}
\begin{tabular}{lllcc}
\hline
Group & Runs & Training paradigm & Max-force [N] & Mean $\|\vec{F}\|$ [N] \\
\hline \hline
A/B/C/E (20 N) & A1, A2, B2, B3, B4, C2, E1, E2 & Joint PPO + estimator & 20 & $\sim 13$ \\
H (50/100 N)   & H1--H19 (50 N and 100 N variants) & Joint PPO + estimator & 50, 100 & $\sim 33,\,66$ \\
S (frozen 50 N) & S1--S9 & Frozen policy + estimator-only & 50 & $\sim 33$ \\
\hline
\end{tabular}
\end{table}

The S-series uses a locomotion checkpoint trained \emph{without} any force
estimate in its observations; the estimator is then trained as a side
channel with the actor frozen. This isolates the estimator architecture
from the policy-adaptation confound that contaminates the H-series.

%==============================================================================
\section{Master results}
\label{sec:eval_master}

\begin{table}[h]
\centering
\caption{Master results. One row per checkpoint, columns are the
chapter's headline metrics. Magnitude-conditional MAE and per-direction
MAE are reported in the per-axis subsections of \cref{sec:eval_findings}.}
\label{tab:eval_master}
\small
\begin{tabular}{l|cccc|cc|c}
\hline
Run & Rel.\,MAE & Ang.\,err. & Rise~time & $\text{MAE}_{[0,5]\text{N}}$ & $\|\hat{\vec{f}}\|_{\text{flat}}$ & $\|\hat{\vec{f}}\|_{\text{slope}}$ & $\|\hat{\vec{f}}\|_{\text{joystick}}$ \\
 & [\%] & [deg] & [ms] & [N] & [N] & [N] & [N] \\
\hline \hline
A1  & -- & -- & -- & -- & -- & -- & -- \\
A2  & -- & -- & -- & -- & -- & -- & -- \\
B2  & -- & -- & -- & -- & -- & -- & -- \\
B3  & -- & -- & -- & -- & -- & -- & -- \\
B4  & -- & -- & -- & -- & -- & -- & -- \\
C2  & -- & -- & -- & -- & -- & -- & -- \\
E1  & -- & -- & -- & -- & -- & -- & -- \\
E2  & -- & -- & -- & -- & -- & -- & -- \\
\hline
H1 (50N)  & -- & -- & -- & -- & -- & -- & -- \\
H1 (100N) & -- & -- & -- & -- & -- & -- & -- \\
H3a (50N) & -- & -- & -- & -- & -- & -- & -- \\
H3a (100N)& -- & -- & -- & -- & -- & -- & -- \\
H9 (50N)  & -- & -- & -- & -- & -- & -- & -- \\
\hline
S1 & -- & -- & -- & -- & -- & -- & -- \\
S2 & -- & -- & -- & -- & -- & -- & -- \\
S3 & -- & -- & -- & -- & -- & -- & -- \\
S4 & -- & -- & -- & -- & -- & -- & -- \\
S5 & -- & -- & -- & -- & -- & -- & -- \\
S6 & -- & -- & -- & -- & -- & -- & -- \\
S7 & -- & -- & -- & -- & -- & -- & -- \\
S8 & -- & -- & -- & -- & -- & -- & -- \\
S9 & -- & -- & -- & -- & -- & -- & -- \\
\hline
\end{tabular}
\end{table}

\begin{figure}[h]
\centering
\includegraphics[width=0.8\textwidth]{images/eval_master_heatmap_placeholder.pdf}
\caption{Heatmap of the master results table. Rows are runs, columns are
metrics, colour is the normalised value (lower is better for all columns
shown). Used as the visual index into the architectural findings below.}
\label{fig:eval_master_heatmap}
\end{figure}

%==============================================================================
\section{Architectural findings}
\label{sec:eval_findings}

Each subsection isolates one architectural axis by contrasting two or more
runs from \cref{tab:eval_master}. The comparison table lists the runs, the
metric most directly affected, and the result. The accompanying figure
visualises the pair.

\subsection{History length}
\label{sec:eval_history}
Runs: A1 ($h=10$) vs.\ A2 ($h=40$). Metric: magnitude-conditional MAE and
responsiveness. \textbf{[result TBD]}.

\begin{figure}[h]
\centering
\includegraphics[width=0.7\textwidth]{images/eval_history_placeholder.pdf}
\caption{Magnitude-conditional MAE for the history-length pair.}
\label{fig:eval_history}
\end{figure}

\subsection{Force dimension}
\label{sec:eval_dim}
Runs: A2 (3D xyz) vs.\ B3 (6D wrench). Metric: angular error and torque-axis
MAE (for B3 only). \textbf{[result TBD]}.

\subsection{Network size}
\label{sec:eval_net}
Runs: B2 (default) vs.\ B3 (big encoder $[256,128]$, head $[64,32]$).
Metric: overall MAE. \textbf{[result TBD]}.

\subsection{Reconstruction loss}
\label{sec:eval_rec}
Runs: A2 (rec on) vs.\ C2 (rec off). Metric: MAE and zero-force behaviour.
\textbf{[result TBD]}.

\subsection{Loss weighting}
\label{sec:eval_loss}
Runs: B3 (baseline) vs.\ B4 (yaw$\times$3 + torque-angle$\times$3). Metric:
yaw-axis angular error, torque-part MAE. \textbf{[result TBD]}.

\subsection{Temporal decay and TCN variants}
\label{sec:eval_tcn}
Runs: S6 (MLP, uniform weights) vs.\ S7 (linear decay) vs.\ S8 (TCN
preprocessor) vs.\ S9 (TCN replacement). Metric: MAE and responsiveness.
\textbf{[result TBD]}.

\subsection{Force-layout choice}
\label{sec:eval_layout}
Runs: S2 (xyz) vs.\ S3 (xy\_yaw). Metric: yaw angular error. \textbf{[result TBD]}.

\subsection{Training-time force distribution}
\label{sec:eval_training_distribution}

This is the central architectural result of the chapter.

During training, forces are sampled per-axis as $F_i\sim\mathcal{U}(0,
F_{\max})$ independently, so the magnitude distribution is the norm of
an i.i.d.\ uniform cube. With $F_{\max}=50\,$N this gives
$\operatorname{mean}\|\vec{F}\|\approx 33\,$N,
$\operatorname{median}\|\vec{F}\|\approx 32\,$N, and
$P(\|\vec{F}\|<10\,\text{N})<0.01$. The estimator therefore sees almost
no small-force examples, independently of how long training runs.

\begin{figure}[h]
\centering
\includegraphics[width=0.7\textwidth]{images/eval_force_histogram_placeholder.pdf}
\caption{Training-time magnitude distribution of
$\|\vec{F}\|$ for $F_{\max}\in\{20,50,100\}\,$N under the current
per-axis-uniform sampler. Low-force bins are empty in all three cases.}
\label{fig:eval_force_histogram}
\end{figure}

The consequence is visible in the magnitude-conditional MAE curve of
every H-series run: accuracy is best near the training-time mean and
degrades below $\sim10\,$N and above the training max. We compare three
checkpoints that differ only in $F_{\max}$:

\begin{table}[h]
\centering
\caption{Magnitude-conditional MAE across training force ranges.
Columns are evaluation bins on $\|\vec{f}^{*}\|$.}
\label{tab:eval_training_distribution}
\begin{tabular}{l|cccc}
\hline
Run & $[0,5]\,$N & $[5,15]\,$N & $[15,30]\,$N & $[30,50]\,$N \\
\hline \hline
A2 (20 N) & -- & -- & -- & -- \\
H1 (50 N) & -- & -- & -- & -- \\
H1 (100 N)& -- & -- & -- & -- \\
\hline
\end{tabular}
\end{table}

\begin{figure}[h]
\centering
\includegraphics[width=0.7\textwidth]{images/eval_binned_mae_placeholder.pdf}
\caption{Magnitude-conditional MAE as a function of applied force
magnitude, one curve per training $F_{\max}$.}
\label{fig:eval_binned_mae}
\end{figure}

\subsection{Auxiliary rewards}
\label{sec:eval_aux}
Runs: (i)~H6 (no est-reward) vs.\ H7 (force-est-accuracy reward);
(ii)~A2 (no compliance reward) vs.\ E1 (compliance-force reward).
Metric: MAE and, for (ii), walker tracking reward under disturbance.
\textbf{[result TBD]}.

\subsection{Frozen versus joint training}
\label{sec:eval_frozen}
Runs: H3a (joint PPO + estimator, 50 N) vs.\ S6 (frozen policy + estimator
only, 50 N). Metric: zero-force behaviour on joystick trajectory
(\cref{sec:eval_self_motion_fps}), and gait-quality delta under no force.
\textbf{[result TBD]}. The frozen-policy setting is the only way to
untangle \emph{estimator} degradation from \emph{policy} degradation,
which is why it exists.

%==============================================================================
\section{Self-motion false positives under joystick driving}
\label{sec:eval_self_motion_fps}

A useful force estimator must report approximately zero when no external
force is applied, even while the robot is walking. The four-segment
joystick trajectory of \cref{sec:eval_protocol_joystick} records
$\|\hat{\vec{f}}\|$ over $4\times\unit[30]{s}$ of undisturbed driving
per checkpoint. We report the per-segment mean and 95th percentile.

\begin{table}[h]
\centering
\caption{Zero-force estimator output under joystick driving, broken out
by walking state. All values are $\|\hat{\vec{f}}\|$ in Newtons.}
\label{tab:eval_self_motion}
\begin{tabular}{l|cc|cc|cc|cc}
\hline
Run & \multicolumn{2}{c|}{Stand} & \multicolumn{2}{c|}{Walk fwd} &
      \multicolumn{2}{c|}{Turn} & \multicolumn{2}{c}{Strafe} \\
    & mean & $P_{95}$ & mean & $P_{95}$ & mean & $P_{95}$ & mean & $P_{95}$ \\
\hline \hline
A2  & -- & -- & -- & -- & -- & -- & -- & -- \\
B3  & -- & -- & -- & -- & -- & -- & -- & -- \\
H3a (50 N) & -- & -- & -- & -- & -- & -- & -- & -- \\
H3a (100 N)& -- & -- & -- & -- & -- & -- & -- & -- \\
S6  & -- & -- & -- & -- & -- & -- & -- & -- \\
\hline
\end{tabular}
\end{table}

\begin{figure}[h]
\centering
\includegraphics[width=0.8\textwidth]{images/eval_joystick_boxplot_placeholder.pdf}
\caption{Distribution of $\|\hat{\vec{f}}\|$ per walking state, one box
per run. Outliers are the dominant failure mode: a single spurious
\unit[15]{N} reading during a fast turn is enough to corrupt a compliant
velocity command.}
\label{fig:eval_joystick_boxplot}
\end{figure}

The expected read-off: (i)~joint-trained H-series runs at 50 N and
especially 100 N exhibit elevated mean and 95th-percentile readings during
fast turns and strafes; (ii)~frozen-policy S-series runs at the same force
level stay closer to zero, isolating the effect to policy adaptation
rather than estimator architecture; (iii)~the slope segment of the
zero-force standing test provides the gravity-bias baseline.

%==============================================================================
\section{Real-world validation}
\label{sec:eval_real}

A subset of the above is reproduced on the physical Go2 robot. The
subset is chosen as the two runs with the lowest simulated relative MAE,
the two runs with the lowest simulated joystick false-positive
$P_{95}$, and one run that is a clear outlier in simulation (as a
negative control).

\begin{table}[h]
\centering
\caption{Subset of checkpoints deployed to the real robot.}
\label{tab:eval_real_runs}
\begin{tabular}{ll}
\hline
Role & Checkpoint \\ \hline \hline
Best sim MAE                     & -- \\
Best sim joystick-FPR            & -- \\
Second-best sim MAE              & -- \\
Second-best sim joystick-FPR     & -- \\
Negative control                 & -- \\
\hline
\end{tabular}
\end{table}

Real-world tests are a strict subset of the protocol:
\cref{sec:eval_protocol} tests 1 (static 360\textdegree{} sweep, with a
calibrated hand pull), 3 (zero-force on a real slope), and 4 (joystick
trajectory). Step-response is skipped for lack of a repeatable impulse
apparatus.

%==============================================================================
\section{Discussion}
\label{sec:eval_discussion}

\subsection{Which design choices were load-bearing}
\textbf{[to be filled from \cref{sec:eval_findings}]}.

\subsection{Training force distribution: what the data says}
\textbf{[to be filled from \cref{sec:eval_training_distribution}]}. The
central claim is that per-axis uniform sampling at the current
$F_{\max}$ does not cover the small-force regime that dominates
deployment, and that either reducing $F_{\max}$ or reshaping the
sampling distribution is a more promising direction than further
tuning the estimator architecture.

\subsection{Deployment implications of the self-motion test}
\textbf{[to be filled from \cref{sec:eval_self_motion_fps}]}. The
working hypothesis is that the self-motion false-positive rate tracks
$F_{\max}$ and training paradigm far more tightly than estimator
architecture.
